% !TeX program = lualatex
% ================================================================
% Indonesian Mathematical Dictionary
%
% Generated by the server using the following settings:
%
%   language            Indonesian (ind)
%   text direction      left to right
%   font                Noto Serif
%   fallback font       Noto Serif
%   built from          latex/dictionary/mathematicalDictionary.tex
%   tier 3 vocabulary   green   RGB 0 128 0
%   English text        black   RGB 0 0 0
%   Indonesian text     purple  RGB 112 48 160
%   layout macros       version 3
%
% Please don't edit any information in this comment block - the server
% compares it against the current settings and rebuilds this file when
% they differ, so changes here would be overwritten. However, feel free
% to make updates to the LaTeX below if you want to edit it.
% ================================================================
%
% ---- What is safe to edit in this file ----------------------------------
%
% Safe:     the Indonesian word and the Indonesian definition - the third
%           and fourth arguments of each entry. That is what this file is for:
%           if a translation reads wrongly to somebody who speaks the language,
%           correct it here and every sheet that uses the word follows.
%
% Not safe: the English word and the English definition - the first and second
%           arguments. The first is how an entry is matched to the shared
%           dictionary; the second is the wording the translation was made
%           from, and is how the server knows the translation is still current.
%           Changing an English definition here does not change the shared
%           dictionary - it only makes this entry look up to date when it is
%           not. Edit the English in the shared dictionary instead, and this
%           file will be brought back into step on the next run.
%
% Entries are added and removed by the server, following the shared dictionary.
% An entry the server cannot read is reported rather than guessed at, so if a
% run complains about this file, look for a missing brace.
% -------------------------------------------------------------------------

\documentclass[11pt]{article}
\usepackage[a4paper,margin=18mm]{geometry}
% Characters this language's font has not got are borrowed from another,
% rather than being silently dropped - see docs/translated-sheet-latex.md
\directlua{%
  if luaotfload and luaotfload.add_fallback then
    luaotfload.add_fallback("eallatinfallback",
      { "Noto Serif:mode=harf;" })
  else
    tex.error("this luaotfload is too old to register a fallback font")
  end
}
\usepackage[english]{babel}
\babelprovide[import]{indonesian}
\babelfont[indonesian]{rm}[Renderer=Harfbuzz, BoldFont={Noto Serif}, ItalicFont={Noto Serif}, BoldItalicFont={Noto Serif}, RawFeature={fallback=eallatinfallback}]{Noto Serif}
\babelfont[indonesian]{sf}[Renderer=Harfbuzz, BoldFont={Noto Serif}, ItalicFont={Noto Serif}, BoldItalicFont={Noto Serif}, RawFeature={fallback=eallatinfallback}]{Noto Serif}
\newcommand{\ealtext}[1]{\foreignlanguage{indonesian}{#1}}
\newcommand{\ealtextblock}[1]{{\raggedright\foreignlanguage{indonesian}{#1}\par}}
% ================================================================
% EAL layout helpers
% ================================================================
\usepackage{xcolor}

\definecolor{ealkeycolour}{RGB}{0,128,0}
\definecolor{ealenglish}{RGB}{0,0,0}
\definecolor{ealtwocolour}{RGB}{112,48,160}

\newsavebox{\ealboxen}
\newsavebox{\ealboxtr}
\newlength{\ealglosswd}
\newlength{\ealglossraise}
\setlength{\ealglossraise}{1.15em}

% a tier 3 word, in English and in translation alike
\newcommand{\ealkey}[1]{\textcolor{ealkeycolour}{#1}}
\newcommand{\ealkeytr}[1]{\textcolor{ealkeycolour}{\ealtext{#1}}}

% One translated word above one English word. Built from kernel
% primitives, never a tabular, which would break wherever the sheet
% loads array, booktabs or tabularx. The box is as wide as the wider of
% the two, so a long translation pushes its neighbours aside instead of
% overlapping them, and it claims its full height so a table row or a
% TikZ node makes room for it.
\newcommand{\ealgloss}[2]{%
  \sbox{\ealboxen}{\ealkey{#1}}%
  \sbox{\ealboxtr}{\small\textcolor{ealtwocolour}{\ealtext{#2}}}%
  \setlength{\ealglosswd}{\wd\ealboxen}%
  \ifdim\wd\ealboxtr>\ealglosswd\setlength{\ealglosswd}{\wd\ealboxtr}\fi
  \leavevmode
  \makebox[\ealglosswd][c]{%
    \makebox[0pt][c]{\usebox{\ealboxen}}%
    \makebox[0pt][c]{\raisebox{\ealglossraise}{\usebox{\ealboxtr}}}%
  }%
}

% opens up line spacing around a block of glosses, so the raised
% translations sit clear of the line above
\newenvironment{ealglossed}{\par\linespread{2.1}\selectfont\ignorespaces}{\par}

% a whole translated sentence above its English counterpart. block level,
% so it does not belong inside a tabular cell
\newcommand{\ealpara}[2]{%
  \par\addvspace{0.5\baselineskip}%
  {\color{ealtwocolour}\ealtextblock{#1}}%
  \penalty200
  {\color{ealenglish}#2\par}%
  \addvspace{0.5\baselineskip}%
}

% ================================================================
% Dictionary layout
% ================================================================
\usepackage{multicol}

\newlength{\ealdictgap}
\setlength{\ealdictgap}{1.2mm}

% english word / english meaning / translated word / translated meaning
\newcommand{\dictentrytr}[4]{%
  \par\addvspace{\ealdictgap}%
  \noindent
  {\bfseries\ealkey{#1}}\quad{\small #2}\par
  \nopagebreak
  {\color{ealtwocolour}\ealtext{{\bfseries #3}~\textemdash{} {\small #4}}\par}%
  \par
}

\pagestyle{empty}
\setlength{\parindent}{0pt}

\begin{document}
{\large\bfseries Indonesian Mathematical Dictionary}\par\medskip

\dictentrytr{coefficient}{the number in front of a letter, showing how many of it there are}{koefisien}{bilangan di depan suatu huruf, yang menunjukkan berapa banyak huruf itu ada}
\dictentrytr{decimal}{a number written with a point, where the digits after the point are parts of a whole}{desimal}{bilangan yang ditulis dengan koma, yang angka-angka setelah koma merupakan bagian dari satu keseluruhan}
\dictentrytr{denominator}{the number below the line in a fraction, showing how many equal parts the whole is split into}{penyebut}{bilangan di bawah garis pada pecahan, yang menunjukkan berapa banyak bagian sama besar keseluruhan itu dibagi}
\dictentrytr{difference}{the answer when one number is subtracted from another}{selisih}{hasil ketika satu bilangan dikurangkan dari bilangan lain}
\dictentrytr{digit}{a single figure from 0 to 9 that a number is written with}{digit}{satu angka dari 0 sampai 9 yang digunakan untuk menulis suatu bilangan}
\dictentrytr{equivalent}{having the same value, even though it is written differently}{senilai}{mempunyai nilai yang sama, walaupun ditulis dengan cara berbeda}
\dictentrytr{estimate}{a sensible answer worked out quickly from rounded numbers}{taksiran}{jawaban masuk akal yang dihitung dengan cepat dari bilangan-bilangan yang dibulatkan}
\dictentrytr{expand}{to multiply out brackets so the expression has no brackets left}{menjabarkan}{mengalikan untuk menghilangkan tanda kurung sehingga ekspresi tidak lagi memiliki tanda kurung}
\dictentrytr{expression}{a group of numbers and letters joined by operations, with no equals sign}{ekspresi}{kelompok bilangan dan huruf yang dihubungkan oleh operasi, tanpa tanda sama dengan}
\dictentrytr{factor}{a whole number that divides exactly into another number}{faktor}{bilangan bulat yang membagi habis bilangan lain}
\dictentrytr{factorise}{to write an expression as a product, usually by putting brackets back in}{memfaktorkan}{menulis suatu ekspresi sebagai hasil kali, biasanya dengan memasukkan kembali tanda kurung}
\dictentrytr{formula}{a rule written with letters, showing how to work one quantity out from others}{rumus}{aturan yang ditulis dengan huruf, yang menunjukkan cara menghitung suatu besaran dari besaran-besaran lain}
\dictentrytr{fraction}{a number written as one whole split into equal parts, with a number of those parts taken}{pecahan}{bilangan yang ditulis sebagai satu keseluruhan yang dibagi menjadi bagian-bagian sama besar, dengan sejumlah bagian itu diambil}
\dictentrytr{gradient}{how steep a line is, found by how far it rises for each step across}{gradien}{seberapa curam suatu garis, yang dicari dari seberapa jauh garis itu naik untuk setiap langkah ke samping}
\dictentrytr{highest common factor}{the largest whole number that divides exactly into two or more numbers}{faktor persekutuan terbesar (FPB)}{bilangan bulat terbesar yang membagi habis dua bilangan atau lebih}
\dictentrytr{identity}{a statement that is true for every value of the letters in it}{identitas}{pernyataan yang benar untuk setiap nilai huruf-huruf di dalamnya}
\dictentrytr{index}{the small raised number showing what power something is raised to}{eksponen}{angka kecil yang ditulis di atas dan menunjukkan pangkat berapa suatu bilangan dipangkatkan}
\dictentrytr{inequality}{a statement that one quantity is larger or smaller than another}{pertidaksamaan}{pernyataan bahwa suatu besaran lebih besar atau lebih kecil daripada besaran lain}
\dictentrytr{integer}{a whole number, which may be positive, negative or zero}{bilangan bulat}{bilangan bulat, yang dapat positif, negatif, atau nol}
\dictentrytr{lowest common multiple}{the smallest number that two or more numbers all divide into exactly}{kelipatan persekutuan terkecil (KPK)}{bilangan terkecil yang habis dibagi oleh dua bilangan atau lebih}
\dictentrytr{multiple}{a number you get by multiplying another number by a whole number}{kelipatan}{bilangan yang diperoleh dengan mengalikan bilangan lain dengan suatu bilangan bulat}
\dictentrytr{multiply}{to add a number to itself a given number of times}{mengalikan}{menambahkan suatu bilangan pada dirinya sendiri sebanyak kali tertentu}
\dictentrytr{numerator}{the number above the line in a fraction, showing how many equal parts are taken}{pembilang}{bilangan di atas garis pada pecahan, yang menunjukkan berapa banyak bagian sama besar yang diambil}
\dictentrytr{percentage}{a number of parts out of a hundred}{persentase}{sejumlah bagian dari seratus}
\dictentrytr{place value}{the value a digit has because of where it sits in a number}{nilai tempat}{nilai yang dimiliki suatu digit karena posisinya dalam suatu bilangan}
\dictentrytr{power}{how many times a number is multiplied by itself}{pangkat}{berapa kali suatu bilangan dikalikan dengan dirinya sendiri}
\dictentrytr{prime}{a whole number greater than 1 that divides only by itself and 1}{prima}{bilangan bulat lebih dari 1 yang hanya habis dibagi oleh dirinya sendiri dan 1}
\dictentrytr{product}{the answer when two or more numbers are multiplied}{hasil kali}{hasil ketika dua bilangan atau lebih dikalikan}
\dictentrytr{proportion}{a relationship where two amounts change by the same factor}{proporsi}{hubungan di mana dua besaran berubah dengan faktor yang sama}
\dictentrytr{quotient}{the answer when one number is divided by another}{hasil bagi}{hasil ketika satu bilangan dibagi oleh bilangan lain}
\dictentrytr{ratio}{a way of comparing two or more amounts, showing how much of one there is for each of the other}{perbandingan}{cara membandingkan dua besaran atau lebih, yang menunjukkan berapa banyak suatu besaran untuk setiap besaran lainnya}
\dictentrytr{reciprocal}{the number you multiply by to get 1, found by turning a fraction upside down}{kebalikan}{bilangan yang kamu kalikan untuk mendapatkan 1, diperoleh dengan membalik suatu pecahan}
\dictentrytr{recur}{to repeat the same digits forever after the decimal point}{berulang}{mengulang angka-angka yang sama terus-menerus setelah koma}
\dictentrytr{root}{a number that gives another number when multiplied by itself a given number of times}{akar}{bilangan yang menghasilkan bilangan lain ketika dikalikan dengan dirinya sendiri sebanyak kali tertentu}
\dictentrytr{round}{to replace a number with a nearby simpler one}{membulatkan}{mengganti suatu bilangan dengan bilangan lain yang lebih sederhana yang dekat dengannya}
\dictentrytr{scale factor}{the number every length is multiplied by when a shape is enlarged}{faktor skala}{bilangan yang digunakan untuk mengalikan setiap panjang ketika suatu bangun diperbesar}
\dictentrytr{simplify}{to write something in its smallest or plainest form without changing its value}{menyederhanakan}{menulis sesuatu dalam bentuk terkecil atau paling sederhana tanpa mengubah nilainya}
\dictentrytr{sum}{the answer when numbers are added together}{hasil penjumlahan}{hasil ketika bilangan-bilangan dijumlahkan}
\dictentrytr{surd}{a root that cannot be written exactly as a fraction, so it is left in root form}{bentuk akar}{akar yang tidak dapat ditulis secara tepat sebagai pecahan, sehingga dibiarkan dalam bentuk akar}
\dictentrytr{unitary}{a method that first finds the value of one part or one item}{satuan}{metode yang pertama-tama mencari nilai satu bagian atau satu barang}

\end{document}
